\documentclass[border=5mm]{standalone}
% ==============================================================================
% SETUP.TEX - CONFIGURACIÓN GLOBAL DEL PROYECTO
% ==============================================================================
% Este archivo centraliza todos los paquetes y estilos.
% Debe incluirse en el preámbulo del libro principal y de los archivos standalone.
% \PassOptionsToPackage{gray}{xcolor}
% --- 1. CODIFICACIÓN E IDIOMA ---
% fix-cm habilita las fuentes Computer Modern en tamaños arbitrarios (por debajo
% de 5 pt, entre otros). Debe cargarse antes que fontenc.
\usepackage{fix-cm}
\usepackage[utf8]{inputenc}
\usepackage[T1]{fontenc}
\usepackage[spanish, es-tabla]{babel}  
\usepackage{csquotes} % Recomendado para babel y citas

\usepackage{import}
% mode=image: Usa el PDF pre-compilado, nunca intenta recompilar.
% Debes compilar cada figura manualmente antes de compilar el libro.
\usepackage[mode=image, subpreambles=false]{standalone}

% --- 2. MATEMÁTICAS Y CIENCIAS ---
\usepackage{amsmath, amssymb, amsfonts, mathtools}
\usepackage{enumitem}

% LaTeX trae \DeclareMathSizes{5}{5}{5}{5}, así que a 5 pt (\tiny) los subíndices
% salen del mismo tamaño que la base: en las etiquetas de figuras el M de
% $\omega_M$ queda desproporcionado. Se restablece la proporción habitual.
\DeclareMathSizes{5}{5}{3.7}{3}

% --- 3. DISEÑO Y GEOMETRÍA --- 
% Unificamos los márgenes para todo el documento
\usepackage[left=2.5cm,right=2.5cm,top=3cm,bottom=3cm]{geometry}
\makeatletter
\ifstandalone % Evita que geometry dañe el recorte de las figuras bajándolas 3cm
  \@ifclasswith{standalone}{crop=false}{
    % Es un capítulo/sección: Dejamos que geometry actúe.
  }{
    % Es una figura: Neutralizamos geometry para que no las recorte mal.
    \geometry{pass}
  }
\fi
\makeatother
\usepackage{parskip} % Espaciado entre párrafos sin sangría
\usepackage{float}   % Para usar [H] en figura

% Ajustes de flotantes para evitar espacios verticales exagerados
\renewcommand{\topfraction}{0.95}      % Permite que las figuras ocupen hasta el 95% superior
\renewcommand{\bottomfraction}{0.95}   % Permite que las figuras ocupen hasta el 95% inferior
\renewcommand{\textfraction}{0.05}     % Permite hojas con tan solo un 5% de texto
\renewcommand{\floatpagefraction}{0.8} % Requiere que una página de solo figuras esté 80% llena
\setcounter{topnumber}{4}
\setcounter{bottomnumber}{4}
\setcounter{totalnumber}{10}

% --- 4. GRÁFICOS Y TABLAS ---
\usepackage{graphicx}
\usepackage{booktabs} % Tablas profesionales
\usepackage{caption}  % Control de captions y \captionof
\usepackage{tikz}
\usetikzlibrary{arrows.meta, calc, angles, quotes, babel, backgrounds}
\usepackage{pgfplots}
\usepgfplotslibrary{groupplots}
\usepgfplotslibrary{external}
\usepgfplotslibrary{patchplots}
\pgfplotsset{compat=1.18}
\usepackage[american, straightvoltages]{circuitikz}

% --- 5. COLORES Y CAJAS ---

\usepackage[table]{xcolor}
\usepackage{tcolorbox}

% Definición de colores corporativos/estilo
\definecolor{utnblue}{RGB}{0, 51, 102}
\definecolor{codegreen}{rgb}{0,0.6,0}
\definecolor{codegray}{rgb}{0.5,0.5,0.5}
\definecolor{codepurple}{rgb}{0.58,0,0.82}
\definecolor{backcolour}{rgb}{0.96,0.96,0.96}

% --- 6. ENTORNOS PERSONALIZADOS (CAJAS) ---

% Caja "Resultado" (Azul Corporativo) — Definiciones, fórmulas clave, resultados importantes.
% Uso: \begin{resultado}[Fórmula de Euler] ... \end{resultado}
\newtcolorbox{resultado}[1][]{
    colback=utnblue!5!white,
    colframe=utnblue,
    title=\textbf{#1},
    fonttitle=\bfseries,
    sharp corners,
    boxrule=0.5mm
}

% Caja "Nota" (Gris) — Advertencias, aplicaciones, contexto secundario.
% Uso: \begin{nota}[Cuidado con la calculadora] ... \end{nota}
% Uso: \begin{nota} ... \end{nota}  → título por defecto "Nota"
\newtcolorbox{nota}[1][Nota]{
    colback=codegray!5!white,
    colframe=codegray,
    title=\textbf{#1},
    fonttitle=\bfseries,
    sharp corners,
    boxrule=0.5mm
}

% Caja inline para resaltar un valor y enlazarlo con su otra aparición en una tabla
% o en una cuenta: mismo color, mismo valor.
% Uso: \marcavalor{$4$}  o  \marcavalor[codegreen]{$4$}
\newtcbox{\marcavalor}[1][utnblue]{
    on line,
    boxsep=1pt, left=2pt, right=2pt, top=1pt, bottom=1pt,
    colback=#1!10!white,
    colframe=#1,
    boxrule=0.4pt,
    arc=2pt
}

% --- 7. CÓDIGO FUENTE (LISTINGS) ---
\usepackage{listings}

\lstdefinestyle{miestilo}{
    backgroundcolor=\color{backcolour},   
    commentstyle=\color{codegreen},
    keywordstyle=\color{magenta},
    numberstyle=\tiny\color{codegray},
    stringstyle=\color{codepurple},
    basicstyle=\ttfamily\footnotesize,
    breakatwhitespace=false,         
    breaklines=true,                 
    captionpos=b,                    
    keepspaces=true,                 
    numbers=left,                    
    numbersep=5pt,                  
    showspaces=false,                
    showstringspaces=false,
    showtabs=false,                  
    tabsize=2,
    frame=single,                    
    rulecolor=\color{codegray}
}

\lstset{style=miestilo}
% Soporte para tildes y ñ en bloques de código
\lstset{literate={ñ}{{\~n}}1 {á}{{\'a}}1 {é}{{\'e}}1 {í}{{\'i}}1 {ó}{{\'o}}1 {ú}{{\'u}}1}

% Operadores matemáticos personalizados

% Vector en sentido discreto (lista finita de coordenadas): negrita + flecha.
% Uso: \vect{v}, \vect{e}_k, \vect{0}.
\newcommand{\vect}[1]{\vec{\mathbf{#1}}}

% Fecha de versión y comando para capítulos standalone
\providecommand{\verdate}{\the\year.\ifnum\month<10 0\fi\the\month.\ifnum\day<10 0\fi\the\day}
\newcommand{\chapterversion}{%
    \vspace{-1.5cm}\begin{flushright}\small\textit{\color{codegray}Versión~\verdate}\end{flushright}\vspace{0.5cm}%
}

% --- 8. HIPERVÍNCULOS (DEBE SER EL ÚLTIMO PAQUETE) ---
\usepackage[colorlinks=true, linkcolor=utnblue, urlcolor=utnblue, citecolor=utnblue]{hyperref}

% --- 9. REFERENCIAS CRUZADAS ENTRE CAPÍTULOS STANDALONE ---
% Cuando se compila un capítulo standalone, xr-hyper lee los .aux de los
% otros capítulos (si existen) para resolver \ref{} a labels externos.
% En la compilación del libro completo este bloque no se activa.
\ifstandalone
  \usepackage{xr-hyper}
  \makeatletter
  \newcommand{\xrchap}[1]{%
    \edef\xrc@self{\detokenize{Capitulo#1}}%
    \edef\xrc@job{\jobname}%
    \ifx\xrc@self\xrc@job\else
      \IfFileExists{../#1capitulo/Capitulo#1.aux}%
        {\externaldocument{../#1capitulo/Capitulo#1}}{}%
    \fi
  }
  \makeatother
  \xrchap{01}\xrchap{02}\xrchap{03}\xrchap{04}
  \xrchap{05}\xrchap{06}\xrchap{07}\xrchap{08}
  \xrchap{09}\xrchap{10}\xrchap{11}\xrchap{12}
  \xrchap{13}
\fi
\usetikzlibrary{calc}

% BIBO: misma entrada x(t) = u(t) acotada por M_x = 1; tres paneles en columna
% con misma escala vertical.
%   Panel 1: entrada x(t) = u(t), cota M_x = 1.
%   Panel 2: salida estable   y_e(t) = 1.5 (1 - e^{-2t}). Parte de 0 y tiende
%            asintoticamente a M_y = 1.5.
%   Panel 3: salida inestable y_i(t) = e^{+t/3} - 1. Parte de 0 y crece sin cota
%            (flecha tangente a la curva + etiqueta "-> inf").

\begin{document}
\begin{tikzpicture}[>=Latex, font=\small]
    \pgfplotsset{
        panel_base/.style={
            width=5.4cm, height=5cm,
            axis lines=middle,
            axis line style={->, thick, black},
            xmin=-0.2, xmax=2.5,
            ymin=-0.45, ymax=2.1,
            xtick=\empty, ytick=\empty,
            ylabel style={font=\small, rotate=0, anchor=south east,
                          at={(axis description cs:0,1)}},
            clip=true,
        },
        bound_line/.style={dashed, gray!55, thin},
        jump_line/.style={dashed, thick},
    }

    % ===== Panel 1: entrada x(t) = u(t) =====
    \begin{axis}[
        name=P1,
        panel_base,
        ylabel={$x(t)$},
    ]
        \addplot[bound_line] coordinates {(-0.18,1) (2.45,1)};
        \node[font=\scriptsize, gray!35!black, anchor=south east]
              at (axis cs: 2.45, 1) {$+M_x$};
        \addplot[thick, utnblue] coordinates {(-0.2,0) (0,0)};
        \addplot[thick, utnblue] coordinates {(0,1) (2.5,1)};
        \addplot[jump_line, utnblue] coordinates {(0,0) (0,1)};
        \node[anchor=north west, inner sep=1pt] at (axis cs: 2.3, -0.07) {$t$};
    \end{axis}
    % leyenda bajo panel 1
    \node[anchor=north, align=center, font=\footnotesize]
        at ($(P1.south)+(0,-0.15cm)$)
        {\textbf{Entrada}\\[1pt]$x(t) = u(t)$};

    % ===== Panel 2: salida estable =====
    \begin{axis}[
        name=P2,
        at={($(P1.east)+(0.4cm,0)$)},
        anchor=west,
        panel_base,
        ylabel={$y_e(t)$},
    ]
        \addplot[bound_line] coordinates {(-0.18,1.5) (2.45,1.5)};
        \node[font=\scriptsize, gray!35!black, anchor=south east]
              at (axis cs: 2.45, 1.5) {$+M_y$};
        \addplot[thick, red!80!black] coordinates {(-0.2,0) (0,0)};
        \addplot[thick, red!80!black, domain=0:2.5, samples=120, smooth]
            {1.5*(1 - exp(-2*x))};
        \node[anchor=north west, inner sep=1pt] at (axis cs: 2.3, -0.07) {$t$};
    \end{axis}
    \node[anchor=north, align=center, font=\footnotesize]
        at ($(P2.south)+(0,-0.15cm)$)
        {\textbf{Salida estable}\\[1pt]$y_e(t) = 1.5\,(1 - e^{-2t})$};

    % ===== Panel 3: salida inestable =====
    % clip=false permite que la flecha de escape y su etiqueta se extiendan
    % por encima de ymax sin recortarse.
    \begin{axis}[
        name=P3,
        at={($(P2.east)+(0.4cm,0)$)},
        anchor=west,
        panel_base,
        clip=false,
        ylabel={$y_i(t)$},
    ]
        \addplot[thick, red!80!black] coordinates {(-0.2,0) (0,0)};
        % signal e^{t/3} - 1 a lo largo del rango visible [0, 2.5].
        % En t=2.5: y = e^{0.833} - 1 ~= 1.300. Pendiente tangente: (1/3)*2.300 ~= 0.767.
        \addplot[thick, red!80!black, domain=0:2.5, samples=100, smooth]
            {exp(x/3) - 1};
        % flecha de escape tangente a la curva (slope ~0.77) + etiqueta
        \draw[->, thick, red!80!black]
              (axis cs: 2.5, 1.30) -- (axis cs: 2.78, 1.515);
        \node[font=\scriptsize, red!80!black, anchor=south west]
              at (axis cs: 2.78, 1.515) {$\to \infty$};
        \node[anchor=north west, inner sep=1pt] at (axis cs: 2.3, -0.07) {$t$};
    \end{axis}
    \node[anchor=north, align=center, font=\footnotesize]
        at ($(P3.south)+(0,-0.15cm)$)
        {\textbf{Salida inestable}\\[1pt]$y_i(t) = e^{+t/3} - 1$};
\end{tikzpicture}
\end{document}
